% ============================================================================
%  chapter13.tex
%  Advanced Memory Management and Custom Allocators
% ============================================================================

\chapter{Advanced Memory Management and Custom Allocators}

\section{The Problem malloc Solves}

Every data structure in Chapter~4 asked \code{malloc} for bytes, and
Chapter~1 taught us to check the result. But \code{malloc} is not a
language feature; it is a \emph{program}---an allocator---that runs
inside your process, and its behavior decides the performance and
fragmentation of everything you build. This chapter opens that black
box: how allocators work, why memory fragments, and how to write your
own allocators when the defaults are the wrong shape for your workload.

The two extremes of allocation, drawn before we go anywhere:

\begin{verbatim}
   From the kernel's point of view (ch.12):          From your program's view:

   the heap is pages the kernel hands you.           the heap is a pool of
                                                     variable-size chunks you
     heap:  [ page ][ page ][ page ][ page ]        request and release:
                                                     malloc(12) -> a 12-byte
   malloc() gets a page when it needs more,          chunk, then free() it.
   free() returns memory to the allocator,           The allocator sits in
   only rarely back to the kernel.                   between, juggling the
                                                     chunks inside the pages.
\end{verbatim}

The allocator's contract is simple and brutal: \emph{given an
arbitrary sequence of malloc and free calls, satisfy every request,
reuse freed memory when possible, waste as little as possible, and do
it fast}. Getting all four right simultaneously is impossible in the
worst case; every allocator is a set of bets about how programs
behave.

\begin{osconcept}
The kernel provides memory only in pages, and only at the granularity
of whole pages. Individual small allocations would be absurdly wasteful
if each were its own page, so the OS delegates the juggling to a
user-space allocator inside your process. That is why \code{malloc}
is a library, not a syscall---it is an OS-adjacent program built on
the kernel's page grants, and its design is where the OS and the
application meet.
\end{osconcept}

\section{Free Lists and the Split/Merge Dance}

The heart of every general allocator is the \emph{free list}: a linked
list of unused chunks, threaded through the freed memory itself. When
you \code{free} a chunk, it is added to the list. When you
\code{malloc}, the allocator searches the list for a chunk big enough,
\emph{splits} it if it is larger than needed, and removes the piece
you get. When a chunk's neighbors are also free, the allocator
\emph{merges} them back into one bigger chunk---the only defense
against fragmentation.

\begin{verbatim}
   The free list, drawn as linked memory:

     heap:  [ A ][ F1 ][ B ][ F2 ][ C ]
                ^      ^
                |      +-- free chunks threaded through the heap

     free list:  F1 --> F2 --> NULL
                 (each free chunk stores the next free chunk's address
                  INSIDE its own space -- free memory is exactly the
                  space you can afford to use for bookkeeping)

   malloc(20) from F1 (size 48):
     - is F1 big enough? yes
     - split:  F1.used = 20 bytes for you; F1.tail = 28 bytes, freed
     - return the 20-byte chunk; the 28-byte remainder rejoins the list

   free(your chunk): add it to the list; if its neighbors are free,
     MERGE them, because 28 + 20 + 16 adjacent beats three scattered
     chunks of any smaller size.
\end{verbatim}

The split/merge cycle is the allocator's entire life. Its difficulty
is that the two operations fight each other: splitting creates small
fragments, and only merging heals them, but merging only works on
\emph{adjacent} freed chunks, and adjacency is not something you can
control after the fact.

\begin{lstlisting}[caption={A minimal free-list allocator, on a static arena.}]
#include <stdio.h>
#include <stdlib.h>
#include <stdint.h>
#include <string.h>

/* The arena: one big byte buffer the allocator hands out from. */
#define ARENA_SIZE 65536
static unsigned char arena[ARENA_SIZE];

/* Each chunk starts with a header describing it. The header lives in
   the arena, right before the user's bytes.                        */
typedef struct chunk {
    size_t size;          /* user-visible size in bytes */
    int free;             /* 1 if free, 0 if allocated */
    struct chunk *next;   /* next chunk in the arena (not the free list) */
} chunk;

#define HEADER_SIZE (sizeof(chunk))

/* The first chunk in the arena. Built lazily on first use. */
static chunk *arena_head = NULL;

static void arena_init(void) {
    chunk *c;

    /* One chunk covering the whole arena, free. */
    c = (chunk *) arena;
    c->size = ARENA_SIZE - HEADER_SIZE;
    c->free = 1;
    c->next = NULL;
    arena_head = c;
}

void *my_malloc(size_t size) {
    chunk *cur;

    if (arena_head == NULL) {
        arena_init();
    }

    /* Round the request up to at least a header, and to alignment. */
    if (size < 16) {
        size = 16;
    }
    size = (size + 7) & ~(size_t) 7;

    /* First-fit search: the first free chunk big enough wins. */
    cur = arena_head;
    while (cur != NULL) {
        if (cur->free && cur->size >= size) {
            break;
        }
        cur = cur->next;
    }
    if (cur == NULL) {
        return NULL;        /* no room */
    }

    /* Split: if the leftover is big enough to be useful, do it. */
    if (cur->size >= size + HEADER_SIZE + 8) {
        chunk *tail;

        tail = (chunk *) ((unsigned char *) (cur + 1) + size);
        tail->size = cur->size - size - HEADER_SIZE;
        tail->free = 1;
        tail->next = cur->next;
        cur->next = tail;

        cur->size = size;
    }

    cur->free = 0;
    return (void *) (cur + 1);
}

void my_free(void *p) {
    chunk *c;

    if (p == NULL) {
        return;
    }
    c = (chunk *) p - 1;     /* the header lives just before the user bytes */
    c->free = 1;

    /* Coalesce with the NEXT chunk if it is free. */
    if (c->next != NULL && c->next->free) {
        c->size = c->size + HEADER_SIZE + c->next->size;
        c->next = c->next->next;
    }
}

void *my_calloc(size_t n, size_t size) {
    void *p;

    p = my_malloc(n * size);
    if (p != NULL) {
        memset(p, 0, n * size);
    }
    return p;
}
\end{lstlisting}

Study the shape of this code, because it is the template for every
allocator in this chapter: a header per chunk, first-fit search,
split on allocation, coalesce-on-next on free. Two details are
deliberate: the header lives \emph{immediately before} the user's
bytes, so \code{my\_free} finds it with \code{(chunk *) p - 1}; and
the coalescing only merges with the \emph{next} chunk, because walking
backward would need a doubly linked structure or a separate metadata
area.

\begin{verbatim}
   The memory layout my_malloc produces:

     +-------+--------------------------+-------+-------------+
     | hdr A |  your 20 bytes           | hdr B | 28 free     |
     +-------+--------------------------+-------+-------------+
       ^ cur+1 returned                  ^ header of the remainder

   free(p) finds header A (p - 1), marks it free, and merges it with
   B if B is free -- so a later malloc(40) can use A+B together.
\end{verbatim}

\begin{osconcept}
The free list is a user-space data structure, but it rests on a kernel
fact: the heap is pages you were granted, and the allocator decides
how to divide them. The kernel cannot merge your fragments, because it
sees only pages, not your chunks. Splitting, merging, and searching
are therefore application-level problems---the OS hands out the clay,
and the allocator sculpts it.
\end{osconcept}

\section{Fragmentation: The Silent Tax}

Why do long-running programs grow forever even when they free
everything? The answer is \emph{fragmentation}, and it comes in two
kinds:

\begin{verbatim}
   INTERNAL fragmentation: you asked for bytes, you got a bigger
   chunk, and the difference is wasted.

     malloc(12)  ->  chunk of 16 (the allocator rounds up):
                     +--------+-------------+
                     | header | 16 bytes    |  4 bytes never used
                     +--------+-------------+
                     internal waste = 4 bytes, invisible to you.

   EXTERNAL fragmentation: free bytes exist, but they are in
   scattered pieces, none big enough for a request.

     heap after a long life:
     [ X ][ F 8 ][ X ][ F 8 ][ X ][ F 8 ][ X ][ F 8 ]
       ^   ^----^   ^      ^      ^      ^
       used  free   used   free   used   free   32 bytes free TOTAL,
       chunks of 8: no single 16-byte request can be satisfied.

   External fragmentation is the killer. It is why a process can have
   100 MB free and still fail a 1 MB malloc.
\end{verbatim}

The classic demonstration is the \emph{allocator stone soup}: allocate
many chunks, free every other one, then try to allocate one big chunk.
The bytes are all there; they are just not adjacent.

\begin{verbatim}
   A, B, C, D allocated.   Free B and D.   malloc(big):

     [A][  F  ][C][  F  ]        free = 2 x 10 bytes, separated
       10 free    10 free        big request of 22: FAILS even though
                                 20 bytes are free, because no single
                                 stretch reaches 22.

   The allocator's only weapon is COALESCING on free (merge adjacent
   neighbors), and even that cannot heal a checkerboard pattern where
   no two free chunks touch.
\end{verbatim}

\textbf{Concept} -- the \emph{binary buddy} allocator attacks external
fragmentation structurally. It divides memory into powers of two:
a request is rounded up to the next power of two, the arena is split
in half recursively to make that size, and frees merge only with their
\emph{buddy}---the adjacent half of the same parent. Because every
chunk has exactly one buddy of its own size, merging is guaranteed to
recover full parent blocks, and the free list becomes a set of
power-of-two lists with \code{O(1)} insertion.

\begin{lstlisting}[caption={The buddy allocator: powers of two, buddy merge.}]
#include <stdio.h>
#include <stdlib.h>

/* Sizes tracked: 2^min .. 2^max, stored as lists of free blocks. */
#define MIN_ORDER 4     /* smallest block: 16 bytes  */
#define MAX_ORDER 12    /* largest block:  4096 bytes */
#define MIN_SIZE  (1 << MIN_ORDER)

typedef struct block {
    struct block *next;
} block;

static block *free_lists[MAX_ORDER + 1];

/* A block's buddy is its neighbor of the same size within the parent. */
static void *buddy_of(void *addr, int order) {
    size_t size;

    size = (size_t) 1 << order;
    return (void *) (((size_t) addr) ^ size);
}

/* Return the largest power-of-two order that fits size. */
static int order_for(size_t size) {
    int order;

    size = size - 1;
    order = MIN_ORDER;
    while ((size >> order) != 0) {
        order = order + 1;
    }
    return order;
}

void *buddy_malloc(size_t size) {
    int order;
    int i;
    block *b;
    void *p;

    order = order_for(size);
    if (order > MAX_ORDER) {
        return NULL;
    }

    /* Find the smallest free list at or above this order. */
    for (i = order; i <= MAX_ORDER; i++) {
        if (free_lists[i] != NULL) {
            break;
        }
    }
    if (i > MAX_ORDER) {
        return NULL;     /* nothing big enough exists */
    }

    /* Take the block and split it downward to the right size. */
    b = free_lists[i];
    free_lists[i] = b->next;

    while (i > order) {
        size_t half;

        half = (size_t) 1 << (i - 1);
        p = (void *) ((size_t) b + half);      /* the buddy half */
        ((block *) p)->next = free_lists[i - 1];
        free_lists[i - 1] = (block *) p;       /* free the buddy */
        i = i - 1;
    }

    return b;
}

void buddy_free(void *addr, int order) {
    void *buddy;
    block *b;

    b = (block *) addr;
    while (order < MAX_ORDER) {
        buddy = buddy_of(addr, order);
        /* If the buddy is on our free list, merge: remove it, double. */
        if (free_lists[order] == buddy) {
            free_lists[order] = free_lists[order]->next;
            if (addr < buddy) {
                b = (block *) addr;
            } else {
                b = (block *) buddy;
            }
            addr = b;
            order = order + 1;
        } else {
            break;
        }
    }

    b->next = free_lists[order];
    free_lists[order] = b;
}
\end{lstlisting}

\begin{verbatim}
   The buddy split, drawn:

     arena of 4096 (order 12):
     [             4096            ]          free_lists[12]
            |  malloc(1500) -> order_for(1500) = 11 (2048)
            v
     [     2048       ][   2048 free   ]      free_lists[11]
       ^ you get this ^ the buddy is freed, can be re-split later

     when you free the 2048, its buddy is checked: if both free, the
     pair merges back into one 4096 -- automatically, because the
     buddy is always exactly adjacent and exactly the same size.
\end{verbatim}

The beauty is the property that makes buddy systems attractive in
operating systems themselves: \emph{merging never has to search}. A
block either has its buddy free (merge, guaranteed adjacent) or not
(split again later). No coalescing pass, no walking. The cost is
internal fragmentation---rounding to powers of two can waste nearly
half a block---which is why real kernels tune the order range and use
other allocators for small objects.

\begin{osconcept}
Fragmentation is the allocator's fundamental adversary, and the kernel
feels its effects at a larger scale: a process that fragments its heap
into holes may still need the OS to grant \emph{new} pages, inflating
its footprint. Buddy systems exist because the kernel itself must
allocate physical pages and merge them back, and it wants merging to
be cheap and structural---powers of two make ``is my neighbor free''
a single XOR. Fragmentation is an OS problem because it is a physical
reality of paged memory, not just a library concern.
\end{osconcept}

\section{Arenas and Object Pools}

The general allocator is flexible and fragmented. The \emph{arena}
(also called a \emph{region} or \emph{pool}) is the opposite bet:
allocate a big block, hand out pieces by bumping a pointer, and free
\emph{everything at once}. It is the right answer when allocations
share a lifetime---a request's scratch space, a connection's
buffers, a frame of a game.

\begin{verbatim}
   Arena allocator: a pointer that only moves forward.

     arena:  [ used ][ used ][ used ][             free          ]
                    ^                          ^
               bump pointer             arena + arena_size

     malloc(n): align n, copy the bump pointer out, advance it.
     free:      NOTHING. There is no individual free.
     reset:     set the bump pointer back to the start; the whole
                arena is free again in one instruction.

   The deal: individual frees are impossible; whole-life resets are
   free. Perfect when every allocation in a phase lives as long as
   the phase.
\end{verbatim}

\begin{lstlisting}[caption={arena.c: an arena allocator.}]
#include <stdio.h>
#include <stdlib.h>
#include <stdint.h>

typedef struct arena {
    unsigned char *mem;      /* the whole block              */
    size_t cap;              /* total capacity in bytes      */
    size_t off;              /* bytes handed out so far      */
} arena;

void arena_init(arena *a, size_t cap) {
    a->mem = malloc(cap);
    if (a->mem == NULL) {
        exit(1);
    }
    a->cap = cap;
    a->off = 0;
}

void *arena_alloc(arena *a, size_t size) {
    void *p;

    /* Align to 8 bytes so any type can live there. */
    size = (size + 7) & ~(size_t) 7;

    if (a->off + size > a->cap) {
        return NULL;         /* arena exhausted */
    }
    p = a->mem + a->off;
    a->off = a->off + size;
    return p;
}

/* Free the ENTIRE arena in one step: rewind the bump pointer. */
void arena_reset(arena *a) {
    a->off = 0;
}

void arena_destroy(arena *a) {
    free(a->mem);
    a->mem = NULL;
    a->cap = 0;
    a->off = 0;
}

int main(void) {
    arena a;
    char *s1;
    char *s2;

    arena_init(&a, 4096);

    s1 = arena_alloc(&a, 32);
    s2 = arena_alloc(&a, 64);
    if (s1 == NULL || s2 == NULL) {
        fprintf(stderr, "arena exhausted\n");
        return 1;
    }

    arena_reset(&a);          /* both s1 and s2 are gone at once */

    arena_destroy(&a);
    return 0;
}
\end{lstlisting}

\begin{verbatim}
   Arena vs malloc at a glance:

                     arena              malloc
   alloc cost        a few ops          search free lists
   free individual   not supported      supported
   free everything   one reset          one call per chunk
   fragmentation     none (linear)      external fragmentation
   best for          same-lifetime      unknown lifetimes,
                     batches            general use
\end{verbatim}

\textbf{Concept} -- an \emph{object pool} is the third classic: instead
of variable sizes, a pool handles objects of \emph{one fixed size}
(threads, connections, nodes). Allocation pops from a free list;
free pushes back. Nothing is ever handed to the OS or fragmented,
because every object is interchangeable and exactly one size.

\begin{lstlisting}[caption={An object pool: fixed-size objects, no fragmentation.}]
#include <stdio.h>
#include <stdlib.h>
#include <stdint.h>

#define POOL_OBJS 64
#define OBJ_SIZE  32

typedef struct pool_obj {
    struct pool_obj *next;   /* free list threaded through free objs */
    unsigned char pad[OBJ_SIZE - sizeof(struct pool_obj *)];
} pool_obj;

typedef struct pool {
    pool_obj *free_head;     /* next free object */
    pool_obj storage[POOL_OBJS];
} pool;

void pool_init(pool *p) {
    size_t i;

    for (i = 0; i < POOL_OBJS; i++) {
        p->storage[i].next = p->free_head;
        p->free_head = &p->storage[i];
    }
}

void *pool_alloc(pool *p) {
    pool_obj *o;

    o = p->free_head;
    if (o == NULL) {
        return NULL;         /* pool exhausted */
    }
    p->free_head = o->next;
    return (void *) o;
}

void pool_free(pool *p, void *ptr) {
    pool_obj *o;

    o = (pool_obj *) ptr;
    o->next = p->free_head;      /* push back on the free list */
    p->free_head = o;
}
\end{lstlisting}

\begin{verbatim}
   The pool's memory, drawn:

     pool.storage:  [o0][o1][o2]...[o63]    one fixed-size array

     free list threads through FREE objects only:
       free_head -> o5 -> o2 -> o31 -> NULL

     alloc: pop head; free: push head. O(1), zero searches, zero
     fragmentation, zero system calls. The pool is the data
     structure of the data structure.
\end{verbatim}

\begin{osconcept}
Arena and pool allocators exploit the kernel's own strategy: hand out
large regions rarely, subdivide them locally, and reclaim wholesale.
The kernel grants pages; the arena bumps a pointer through them; the
pool never asks the kernel again. Choosing the allocator is choosing
how often your program crosses the kernel boundary---and good
allocator design is simply the discipline of crossing it as rarely as
possible.
\end{osconcept}

\section{Practical: A Connection Pool for a Server}

Let us apply the chapter to the book's own recurring hero---the
server. A connection pool hands out fixed-size connection objects,
reuses them across clients, and never fragments and never calls
\code{malloc} after startup. This is the object pool of the previous
section wearing its production costume.

\begin{lstlisting}[caption={connpool.c: fixed-size connection objects, recycled.}]
#include <stdio.h>
#include <stdlib.h>

typedef struct conn {
    int fd;                /* the socket (0 = free slot)   */
    int in_use;            /* 1 while a client owns it     */
    struct conn *next;     /* free-list threading           */
} conn;

typedef struct conn_pool {
    conn *slots;           /* the fixed array of conns      */
    int capacity;
    conn *free_head;       /* head of the free list         */
} conn_pool;

int conn_pool_init(conn_pool *p, int capacity) {
    int i;

    p->slots = malloc(sizeof(conn) * (size_t) capacity);
    if (p->slots == NULL) {
        return -1;
    }
    p->capacity = capacity;
    p->free_head = NULL;

    /* Thread every slot onto the free list. */
    for (i = 0; i < capacity; i++) {
        p->slots[i].fd = 0;
        p->slots[i].in_use = 0;
        p->slots[i].next = p->free_head;
        p->free_head = &p->slots[i];
    }
    return 0;
}

/* Hand out a free connection object, or NULL if the pool is full. */
conn *conn_pool_acquire(conn_pool *p) {
    conn *c;

    c = p->free_head;
    if (c == NULL) {
        return NULL;
    }
    p->free_head = c->next;
    c->in_use = 1;
    return c;
}

/* Return a connection to the pool. */
void conn_pool_release(conn_pool *p, conn *c) {
    c->in_use = 0;
    c->fd = 0;
    c->next = p->free_head;
    p->free_head = c;
}

int main(void) {
    conn_pool pool;
    conn *c1;
    conn *c2;

    if (conn_pool_init(&pool, 8) < 0) {
        return 1;
    }

    c1 = conn_pool_acquire(&pool);
    c2 = conn_pool_acquire(&pool);
    if (c1 == NULL || c2 == NULL) {
        fprintf(stderr, "pool exhausted\n");
        return 1;
    }
    printf("acquired two connections\n");

    conn_pool_release(&pool, c1);
    conn_pool_release(&pool, c2);

    /* Reacquire the same objects -- same addresses, no malloc. */
    c1 = conn_pool_acquire(&pool);
    printf("reacquired (pool recycles storage)\n");
    conn_pool_release(&pool, c1);

    free(pool.slots);
    return 0;
}
\end{lstlisting}

\begin{verbatim}
   Why a pool and not malloc per connection:

     malloc/free per connection   pool
     -------------------------    ----
     searches free lists          O(1) pop/push
     can fragment over a day      zero fragmentation (fixed size)
     may cross the kernel         no syscalls after init
     allocator lock contention    single-threaded pop/push (or one
                                   lock for the pool)
     allocator can't recycle      addresses recycle: caches stay warm

   The price: a fixed maximum (capacity) and objects that cannot
   outlive the pool. Both are features when you know your limits.
\end{verbatim}

\begin{researchnote}[The measured case for custom allocators]
Chapter~1 introduced the empirical study by Berger, Zorn, and
McKinley that most custom allocators do not beat \code{malloc}. Now
that you can read their argument in context, the nuance matters: they
found that allocators win when they \emph{specialize to a workload's
lifetime patterns}---exactly what arenas and pools do here---and lose
when they merely reimplement a general allocator slightly differently.

\begin{verbatim}
   The paper's real conclusion, restated honestly:

       "Generic custom allocators rarely help; allocators that match
        a program's allocation AND deallocation patterns -- regions,
        pools, per-thread caches -- are the ones that show gains."

   This chapter is that sentence in practice: the arena (reset in
   O(1)), the pool (fixed size, no fragmentation), and the buddy
   (structural merge) are lifetime-aware designs. A hand-rolled
   malloc-lookalike is the one approach the research says to skip.
\end{verbatim}

So the research note from Chapter~1 and this chapter agree: write
\emph{specialized} allocators when your workload has structure, not
generic ones on a whim. Measure first; the paper's methodology (link
in many allocators, time them) is the correct one to steal.

\paperref{E.~D.~Berger, B.~G.~Zorn, K.~S.~McKinley, ``Reconsidering
Custom Memory Allocation,'' ACM OOPSLA 2002}{https://people.cs.umass.edu/~emery/pubs/berger-oopsla2002.pdf}
\end{researchnote}

\section{Structures in Bytes}

Every allocator in this chapter hides a header in front of the bytes
you receive. On a 64-bit machine the structures are:

\begin{verbatim}
   my_malloc chunk header  24 bytes   offset size field
                                      0      8    size   (user size)
                                      8      4    free
                                      12     4    pad
                                      16     8    next   (next in arena)
   arena                   24 bytes   0      8    mem
                                      8      8    cap
                                      16     8    off
   pool_obj                32 bytes   0      8    next   (free-list link)
                                      8      24   pad[24]
   buddy free block         8 bytes   0      8    next   (threaded inside)
\end{verbatim}

The header sits immediately before the user's bytes, so
\code{my\_free} recovers it with \code{(chunk *) p - 1}; because the
header is 24 bytes and the sizes are rounded to 8, every user pointer
is 8-byte aligned. The arena is a bump pointer plus a capacity, 24
bytes total; the pool is one array of fixed-size objects whose free
list is threaded through the objects themselves. These byte counts are
the whole point: an allocator is a small set of byte records and the
rules for moving them.

\section{Summary}

\begin{verbatim}
   The allocator toolbox, complete:

   free list        freed chunks threaded through memory
   first-fit        the simplest search; split on alloc
   coalescing       merge adjacent frees (the anti-fragmentation cure)
   internal frag    waste inside your own chunk (rounding)
   external frag    free bytes in pieces too small to use
   buddy system     powers of two; buddy merge is O(1)
   arena            bump pointer, reset everything at once
   object pool      one fixed size, O(1) pop/push, no frag

   Rules that keep allocator code correct:
     1. every alloc returns NULL or a valid pointer -- check it;
     2. a header is reachable from the user pointer (p - 1);
     3. coalesce with the NEXT chunk on free, at minimum;
     4. match the allocator to the LIFETIME pattern, not to a mood;
     5. measure before you write one; the default is usually fine.
\end{verbatim}

\begin{osconcept}
Allocators are where the application and the OS shake hands: the
kernel grants pages, the allocator divides them, and fragmentation
wanders back up the chain as wasted pages. Every allocator in this
chapter is a different answer to the same kernel-shaped question--
how rarely can I ask for more?--and the best ones ask almost never.
This is the deepest layer of memory management before we leave the
heap for good.
\end{osconcept}

With memory mastered, the next chapter switches from bytes to
\emph{symbols}: parsing text with state machines, the technique that
powers every config parser, every protocol decoder, and every
compiler's front end.

